\exerciseeditionsection{Solution to Wald's Problem}
\ExerciseEditorialNote

\begin{enumerate}[leftmargin=2.8em,itemsep=1em]
\WaldSolution{1}
Let the garage be at rest in coordinates \((t,x)\), with its door at
\(x=0\) and back wall at \(x=L\).  Let the car move with speed \(v\), and
let its proper length also be \(L\).  Its length on a garage-time slice is
\(L/\gamma\), where \(\gamma=(1-v^2)^{-1/2}\).  When the rear reaches the
door, the front is therefore at \(x=L/\gamma<L\).  On that simultaneity
slice the whole car lies inside the garage, so the doorman's statement is
correct.

In the car frame the garage has length \(L/\gamma\).  The event at which the
front hits the back wall occurs before, in car-frame time, the event at which
the rear reaches the door.  Indeed, if the latter event is the origin, then
the collision event has garage coordinates
\[
 (t,x)=\left(-\frac{L-L/\gamma}{v},L\right),
\]
and hence car time
\[
 t'=\gamma(t-vx)<0.
\]
Thus the driver says that the front has already broken through the wall when
the rear reaches the door.  That statement is also correct.  The two accounts
compare the collision with different distant events because the hypersurfaces
\(t=\text{constant}\) and \(t'=\text{constant}\) are tilted relative to one
another.  The spacetime diagram contains four straight world lines, two
vertical garage boundaries and two parallel timelike car boundaries; the two
families of simultaneity lines display the resolution directly.
\end{enumerate}
